%% ---------------------------------------------------------------------
\section{Espacios métricos: medir distancias}
\label{sec:pm-metricos}
%% ---------------------------------------------------------------------

Tres conjuntos muy distintos van a acompañarnos en esta sección:
\begin{itemize}
  \item $\R^n$, cuyos elementos son vectores $x=[x_1,\dots,x_n]\T$. En un sistema dinámico, el \emph{estado en un instante}.
  \item $C[a,b]$, el conjunto de funciones continuas $f:[a,b]\to\R$. En un sistema dinámico, una \emph{trayectoria completa} o una señal.
  \item Un conjunto de palabras de cuatro letras, $X=\{\text{casa},\text{cama},\text{masa},\text{pasa},\dots\}$.
\end{itemize}
¿Qué tan lejos está un vector de otro? ¿Una función de otra? ¿Una palabra de otra? La respuesta no viene con los objetos: depende de la estructura que se les agregue.

\begin{definicion}[label=def:pm-metrica]{Métrica y espacio métrico}
Sea $X$ un conjunto. Una función $d:X\times X\to\R$ es una \clave{métrica} si para todo $x,y,z\in X$:
\begin{enumerate}[label=(M\arabic*),leftmargin=3em]
  \item $d(x,y)\ge0$ \hfill(no negatividad)
  \item $d(x,y)=0$ si y solo si $x=y$ \hfill(separa puntos)
  \item $d(x,y)=d(y,x)$ \hfill(simetría)
  \item $d(x,y)\le d(x,z)+d(z,y)$ \hfill(desigualdad triangular)
\end{enumerate}
El par $(X,d)$ es un \clave{espacio métrico}.
\end{definicion}

La desigualdad triangular es la que hace el trabajo: dice que no hay atajos, que pasar por un tercer punto nunca acorta el camino. Es la única de las cuatro que hay que probar con cuidado en cada ejemplo.

\begin{ejemplo}[label=ej:pm-metricasRn]{Tres distancias entre los mismos dos puntos}
En $\R^n$ se usan
\[
  d_1(x,y)=\sum_{i=1}^n\abs{x_i-y_i},\qquad
  d_2(x,y)=\Bigl(\sum_{i=1}^n\abs{x_i-y_i}^2\Bigr)^{1/2},\qquad
  d_\infty(x,y)=\max_i\abs{x_i-y_i},
\]
y, en general, $d_p(x,y)=\bigl(\sum_i\abs{x_i-y_i}^p\bigr)^{1/p}$ con $1\le p<\infty$ ($d_\infty$ es el límite de $d_p$ cuando $p\to\infty$). Calcule las tres para $x=[1,1]\T$ y $y=[4,3]\T$, y pruebe que $d_1$ es una métrica.
\solucion $d_1=3+2=5$ (el taxista que solo se mueve por calles), $d_2=\sqrt{9+4}=\sqrt{13}\approx3.61$ (la línea recta) y $d_\infty=\max\{3,2\}=3$ (la coordenada que más difiere). Los puntos son los mismos; la distancia cambia con la métrica (\cref{fig:pm-distancias}, izquierda). Para $d_1$, (M1)--(M3) son inmediatas. Para (M4), en cada componente $x_i-y_i=(x_i-z_i)+(z_i-y_i)$ y la desigualdad triangular de los reales da $\abs{x_i-y_i}\le\abs{x_i-z_i}+\abs{z_i-y_i}$; sumando sobre $i$ se obtiene $d_1(x,y)\le d_1(x,z)+d_1(z,y)$.
\end{ejemplo}

\begin{ejemplo}[label=ej:pm-hamming]{Una distancia entre palabras}
En el conjunto de palabras de cuatro letras, sea $d(u,v)$ el número de posiciones en las que $u$ y $v$ difieren (distancia de Hamming). Calcule $d(\text{casa},\text{cama})$, $d(\text{casa},\text{masa})$ y $d(\text{cama},\text{masa})$, y diga por qué es una métrica.
\solucion $d(\text{casa},\text{cama})=1$, $d(\text{casa},\text{masa})=1$ y $d(\text{cama},\text{masa})=2$. Es una métrica: (M1)--(M3) son claras, y para (M4), si $u$ y $v$ difieren en la posición $i$, entonces $w$ difiere de $u$ o de $v$ en esa posición; cada posición que cuenta en $d(u,v)$ cuenta al menos una vez en $d(u,w)+d(w,v)$. La misma distancia mide cuántos bits cambió el ruido en un mensaje digital: los objetos no tienen que ser números para que tenga sentido hablar de cercanía.
\end{ejemplo}

\script{cap03_preliminares/C3_Ej1_Metricas.m}
\begin{figure}[ht]
\centering
\includegraphics[width=\linewidth]{cap03/f_distancias}
\caption{Izquierda: tres métricas en $\R^2$ para los mismos dos puntos. Derecha: entre funciones de $C[0,1]$, $d_1$ es el área entre las curvas y $d_\infty$ la mayor separación vertical.}
\label{fig:pm-distancias}
\end{figure}

Entre funciones continuas $f,g\in C[a,b]$ se definen, de la misma manera,
\begin{equation}
  d_p(f,g)=\Bigl(\int_a^b\abs{f(\tau)-g(\tau)}^p\,\dif\tau\Bigr)^{1/p},\quad 1\le p<\infty,\qquad
  d_\infty(f,g)=\max_{a\le t\le b}\abs{f(t)-g(t)} .
  \label{eq:pm-dfunciones}
\end{equation}
La suma sobre las componentes se reemplaza por una integral sobre el tiempo. Para $f(t)=\sin2\pi t$ y $g(t)=t$ en $[0,1]$, $d_1(f,g)=0.921$ y $d_\infty(f,g)=1.763$, alcanzada en $t=0.775$ (\cref{fig:pm-distancias}, derecha). Que $d_p(f,g)=0$ implique $f=g$ usa la continuidad: una función continua y no negativa con integral nula es idénticamente cero.

\codigorepo{cap03_preliminares/C3_Ej1_Metricas.m}{vectores}{Tres métricas para los mismos dos puntos}

\begin{trampa}[label=tr:pm-parece]{Parecer una distancia no basta}
\plausible{$d(x,y)=(x-y)^2$ es no negativa, simétrica y vale cero solo si $x=y$; por lo tanto es una métrica en $\R$.}
Falta la desigualdad triangular, y falla: $d(0,2)=4$, pero $d(0,1)+d(1,2)=1+1=2$. Pasar por $z=1$ ``acorta'' el camino. Elevar al cuadrado castiga las distancias grandes más que proporcionalmente, y eso es incompatible con (M4).
\correcto{Para concluir que una función es una métrica hay que verificar las cuatro propiedades; la triangular es la que casi siempre decide. En cambio, $\abs{x-y}^{1/2}$ sí es una métrica en $\R$ (pruébelo: $\sqrt{a+b}\le\sqrt a+\sqrt b$).}
\end{trampa}

\codigorepo{cap03_preliminares/C3_Ej1_Metricas.m}{nometrica}{$(x-y)^2$ no es métrica}

\begin{pausa}
En un tablero de ajedrez, ¿qué métrica cuenta el número de movidas que necesita un rey para ir de una casilla a otra? ¿Y una torre que solo puede moverse una casilla a la vez?
\respuesta{El rey se mueve una casilla en cualquier dirección, incluidas las diagonales: necesita $\max\{\abs{\Delta x},\abs{\Delta y}\}$ movidas, la métrica $d_\infty$. La torre restringida a una casilla por movida necesita $\abs{\Delta x}+\abs{\Delta y}$, la métrica $d_1$.}
\end{pausa}
